\subsection{An exact theta witness on the numerical-obstruction curve}
\label{sec:a27-theta-seven}

\paragraph{Finding A31 (a degree-seven geometric witness on A27).}
Let $k=\mathbf F_4=\mathbf F_2(u)$, with $u^2+u+1=0$, and let
\[
 C/k:\qquad y^2+y=u t^3+t+u/t
 \tag{TS.1}
\]
be the smooth proper rank-one genus-two curve of A27.  Let $W$ be
its point over $t=0$ and let $R=(u,u^2)$.  Then
$M=\mathcal O_C(R-W)$ has exact order seven.  On the Frobenius
twist $C^{(1)}$, its twist $M^{(1)}$ lies on the Raynaud theta
divisor.  Consequently A11 gives a geometrically connected cyclic
degree-seven cover with geometric primitive $2$-rank zero.  After
a finite constant extension it gives a constant split witness at
any prescribed point of $C$.

For the prescribed point $R$, the canonical class
$\kappa_R=[K_C-2R]$ has exact order seven.  Thus A29's detector
does not apply to this class, even after extending constants,
because $\operatorname{ord}_7(2)=3$ is odd.  This is a concrete
instance of the remaining detector case that the geometric method
does handle.  It is not a universal incidence theorem.  It is also
compatible with every numerical failure proved in A27.

\begin{proof}
\emph{1. An explicit isomorphic model.}
The coordinate change
\[
 \begin{aligned}
 X&=u/t, &
 Y&=u^2t^{-2}(y+u)+u^2,\\
 t&=u/X, &
 y&=u+(Y+u^2)/X^2
 \end{aligned}
 \tag{TS.2}
\]
and its displayed inverse identify the function field of (TS.1)
with that of
\[
 \begin{aligned}
 C':\quad Y^2+X^2Y&=f(X),\\
 f(X)&=X^5+X^4+uX^3+u^2X^2+uX+u.
 \end{aligned}
 \tag{TS.3}
\]
Here is a coefficient check, retaining all constant terms.
Put $\alpha=u^2t^{-2}$ and $\beta=t^{-2}+u^2$, so
$Y=\alpha y+\beta$.  Since $X^2\alpha=\alpha^2$,
the equation on the left becomes
\[
 \begin{aligned}
 Y^2+X^2Y
 &=\alpha^2(ut^3+t+u/t)+\beta^2+X^2\beta\\
 &=u^2t^{-5}+ut^{-4}+ut^{-3}
       +ut^{-2}+u^2t^{-1}+u\\
 &=f(u/t).
 \end{aligned}
 \tag{TS.4}
\]
Both are models of their smooth proper function-field curves, so
this birational identification extends to an isomorphism.
The two odd Artin--Schreier pole orders in (TS.1) are three and
one; thus the genus and $2$-rank are those already checked in A27.
The point $W$ becomes the unique point at infinity in (TS.3),
with $\operatorname{ord}_W(X)=-2$ and
$\operatorname{ord}_W(Y)=-5$.  The point $R$ becomes $(1,u)$.
The affine model is smooth: the only possible singular
$X$-coordinate is zero, and $f'(0)=u\ne0$ there.

\emph{2. A principal divisor of degree seven.}
Set
\[
                F=(X+u)Y+X.
 \tag{TS.5}
\]
The hyperelliptic involution sends $Y$ to $Y+X^2$.
Taking its function-field norm gives the exact polynomial identity
\[
 \begin{aligned}
 F\,\iota(F)
 &=X^2+X^3(X+u)+(X+u)^2f(X)\\
 &=(X+1)^7.
 \end{aligned}
 \tag{TS.6}
\]
Indeed expansion yields coefficient one in every degree from
zero through seven.  At the two points above $X=1$,
\[
       F(1,u)=0,\qquad F(1,u^2)=u^2\ne0.
 \tag{TS.7}
\]
They are distinct unramified points since $X^2=1$ there.
Thus (TS.6) assigns all seven zeros of $F$ to $R$, with no
zero at its conjugate.  The function is regular on the affine
curve and has a pole of exact order seven at $W$:
the term $XY$ has pole order seven, whereas $uY$ and $X$
have pole orders five and two.  Consequently
\[
                       \operatorname{div}(F)=7(R-W).
 \tag{TS.8}
\]
The class $[R-W]$ is nonzero.  Otherwise a function with divisor
$R-W$ would give a degree-one map from a genus-two curve to
$\mathbf P^1$, which is impossible.  Since seven is prime,
(TS.8) proves that its exact order is seven.

\emph{3. The theta bundle and the Frobenius twist.}
The differential
\[
                    \omega=dX/X^2=d(1/X)
 \tag{TS.9}
\]
is exact and has divisor $2W$.  At a finite point with $X\ne0$,
$X$ is an etale local coordinate, so $\omega$ is nonvanishing.
At $X=0$, the differentiated equation gives
$\omega=dY/f'(X)$; since $f'(0)=u\ne0$, this is again a
nonvanishing regular differential.  Its divisor is therefore
supported at $W$ and has degree $2g-2=2$.

Let $B_C$ denote the line of locally exact differentials on
$C^{(1)}$.  The section supplied by (TS.9) has divisor
$W^{(1)}$ on that twist: a local zero of order two in an exact
differential corresponds to a zero of order one in the
$\mathcal O_{C^{(1)}}$-line.  Hence
\[
 B_C\simeq\mathcal O_{C^{(1)}}(W^{(1)}),\qquad
 B_C\otimes M^{(1)}
       \simeq\mathcal O_{C^{(1)}}(R^{(1)}).
 \tag{TS.10}
\]
The second line bundle has a nonzero section, proving the asserted
theta incidence.  This keeps the twist explicit over $\mathbf F_4$:
the coefficients of (TS.3) are squared, and
$R^{(1)}=(1,u^2)$ in the resulting coordinates.
Twisting (TS.8) preserves the exact order seven.

\emph{4. The cover and the exact limitation of the detectors.}
The powers of two modulo seven are $1,2,4$.  Adjoining their
negatives gives every nonzero residue, so
\[
                  \langle2,-1\rangle=(\mathbf Z/7)^\times.
 \tag{TS.11}
\]
Finding A11 applies to (TS.10) and gives a cyclic degree-seven
cover with primitive geometric $2$-rank zero.  Since the degree
is prime, this is the whole new factor.  Descending the cover and
its deck transformations and rationalizing a point above any
prescribed base point require only a finite constant extension.
The geometric rank-zero property survives these operations; the
relative rational class number is then odd, and A11's pointed
deck argument separates any putative collision.

Finally (TS.9) also gives $K_C\sim2W$, so
$\kappa_R=-2[R-W]$ has exact order seven.  Its only nonzero odd
primary part is the one for seven, whose order of two is odd.
Injectivity of rational Jacobian points under constant extension
preserves that order and prime support.  Thus this prescribed
point remains outside A29's sufficient hypothesis.  In fact A27
gives $\#J_C(\mathbf F_4)=28$ and $\#C(\mathbf F_4)=6$, so for
every one of those six prescribed points the odd primary support
of $\kappa_P$ is contained in $\{7\}$.  A29's detector therefore
fails for all six points, and continues to fail for these same
points after every constant extension, whereas the geometric
cover just constructed becomes split at each prescribed point
after a finite extension.  A27 refuted
universal success of its stated size inequalities at primes with
two primitive; seven does not satisfy that primitive-root
hypothesis.  Neither A27 nor the failure of A29's detector
asserted absence of the present geometric witness.
\end{proof}

The model (TS.3), norm (TS.6) and isomorphism (TS.2) were found
in this run by the root investigator.  The compact certificate
\texttt{verify\_A31\_norm7.py} checks their coefficients directly
over $\mathbf F_4$, including the inverse substitution, both point
values, smoothness at $X=0$ and the signed unit orbit.  It uses no
Jacobian arithmetic package.  The divisor, differential, exact-order
and covering arguments are the written proof above, not an
inference from a numerical search alone.
